(1) 轨道 \(\phi_\zeta\) 的归一性由 \(4\pi(\zeta^3/\pi)\int_0^\infty r^2e^{-2\zeta r}dr=1\) 得到，乘积波函数也归一且在交换两电子时对称。费米子总态需反对称，所以空间对称试探对应自旋反对称的单态；双电子同轨道的自旋三态违反 Pauli 原理。
(2) 对每个指数轨道，动能期望 \(\langle-\nabla^2/2\rangle=\zeta^2/2\)，与核的Coulomb势能期望 \(\langle-Z/r\rangle=-Z\zeta\)。两电子合计为 \(\zeta^2-2Z\zeta\)；电子间直接排斥积分 \(5\zeta/8\) 为正，所以 \(E(\zeta)=\zeta^2-2Z\zeta+5\zeta/8\)（单位为 Hartree）。
(3) 求导 \(dE/d\zeta=2\zeta-2Z+5/8\)，极小点 \(\zeta_*=Z-5/16\)，二阶导恒为 \(2>0\)。当 \(Z>5/16\) 时此点合法，最小试探能量为 \(E_{\rm var}=-(Z-5/16)^2\) Hartree；氦核 \(Z=2\) 时 \(\zeta_*=27/16\)，\(E_{\rm var}=-(27/16)^2\) Hartree。
(4) Rayleigh–Ritz 原理给真实基态 \(E_0\le\langle\Psi_\zeta|H|\Psi_\zeta\rangle\)，故 \(E_{\rm var}\) 是上界而非精确本征值。\(\zeta_*=Z-5/16<Z\) 反映其他电子屏蔽核电荷，简单独立轨道乘积无法完整描述两电子瞬时排斥关联；若 \(Z\le5/16\)，形式驻点非正，超出此单参数束缚试探的有效极小域，不能使用该表达式宣布存在真正束缚态。
